\section{Surprisal moments of independent copies}
\label{sec:schur-label-iid-moments}

The independent-copy step in
\cite[Section~7, lines~75--85 and 227--238]{OpenAI2026AreaLaw}
uses only the positive eigenvalues of a density matrix. We retain this
convention by weighting the exponential by the eigenvalue itself.

\begin{lemma}[Multiplicative functions of tensor powers]
\label{lem:cfc-finite-tensor-power}
\lean{Matrix.finKronecker_cons,Matrix.cfc_finKronecker_const}
\leanok
Let $A\ge0$ act on a finite-dimensional Hilbert space. Suppose
$f:\mathbb R\to\mathbb R$ satisfies $f(1)=1$ and
$f(xy)=f(x)f(y)$ for $x,y\ge0$. Then, for every integer $k\ge0$,
\[
 f(A^{\otimes k})=f(A)^{\otimes k}.
\]
Here the zero-copy tensor power is the identity on $\mathbb C$.
More generally, separating the first factor of a finite tensor product
agrees with the canonical identification of a configuration with its
first coordinate and its remaining coordinates.
\end{lemma}

\begin{proof}
\leanok
Separate the first factor and apply the multiplicative functional calculus
for two positive matrices. Induction proves the assertion; the zero-copy
case is precisely $f(1)=1$. Since the spectra are finite, no continuity
hypothesis on $f$ is needed.
\end{proof}

\begin{lemma}[Exact centered surprisal moments]
\label{lem:schur-label-iid-moment-identity}
\lean{Matrix.re_trace_mul_exp_centered_neg_log_eq_surprisalMoment,
Matrix.re_trace_mul_exp_centered_neg_log_finKronecker}
\uses{lem:cfc-finite-tensor-power}
\leanok
For a Hermitian matrix $A$ with eigenvalues $\lambda_j$, real $s,t$, and
$L_A=-\log A$ defined by spectral calculus, one has
\[
 \operatorname{ReTr}\bigl(Ae^{t(L_A-sI)}\bigr)
 =e^{-ts}\sum_j\lambda_j e^{-t\log\lambda_j}.
\]
The real logarithm here is defined at zero by zero and, for a negative
argument, as the logarithm of its absolute value. In particular, for
$A\ge0$, each zero eigenvalue contributes zero, including when $t=1$.
If $\rho_k=A^{\otimes k}$, then
\[
 \operatorname{ReTr}\bigl(\rho_k e^{t(L_{\rho_k}-ksI)}\bigr)
 =\left[\operatorname{ReTr}\bigl(Ae^{t(L_A-sI)}\bigr)\right]^k.
\]
This last assertion holds for all real $s,t$ and every $k\ge0$.
\end{lemma}

\begin{proof}
\leanok
The first equality follows by diagonalization. For the second, apply the
preceding lemma to $f_t(x)=x e^{-t\log x}$. This function vanishes at zero;
when both factors are positive, its multiplicativity follows from
$\log(xy)=\log x+\log y$. The trace of a tensor product is the product
of the traces, and $e^{-tks}=(e^{-ts})^k$.
\end{proof}

\begin{theorem}[Label moments from the one-copy hypothesis]
\label{thm:schur-label-iid-mgf}
\lean{TensorPower.log_centered_surprisalMoment_finKronecker_le,
TensorPower.log_centered_labelEntropy_finKronecker_le}
\uses{lem:schur-label-iid-moment-identity,thm:schur-label-moment-bounds}
\leanok
Let $\sigma\ge0$ have trace one on a space of dimension $q$, and write
$S(\sigma)=-\sum_j\lambda_j\log\lambda_j$. Suppose $C_0\ge0$,
$c_0>0$, $\mathcal B\ge1$, and
\[
 \log\sum_j\lambda_j e^{-t\log\lambda_j}
 \le tS(\sigma)+C_0\mathcal B t^2
 \qquad\left(|t|\le\frac{c_0}{\sqrt{\mathcal B}}\right).
\]
For $\rho_k=\sigma^{\otimes k}$ and $L_k=-\log\rho_k$ on its support,
\[
 \log\operatorname{ReTr}\bigl(\rho_k e^{t(L_k-kS(\sigma)I)}\bigr)
 \le C_0k\mathcal B t^2
 \qquad\left(|t|\le\frac{c_0}{\sqrt{\mathcal B}}\right).
\]
Let $F_k$ be the actual Schur-label observable for the permutation of these
$k$ copies, and put $c=\min(c_0,1)/2>0$. Then
\[
 \log\operatorname{ReTr}\bigl(\rho_k e^{t(F_k-kS(\sigma)I)}\bigr)
 \le 2C_0k\mathcal B t^2+\frac{q^2}{2}\log(k+1)
 \qquad\left(|t|\le\frac{c}{\sqrt{\mathcal B}}\right).
\]
No full-rank hypothesis is required.
\end{theorem}

\begin{proof}
\leanok
The scalar one-copy moment is strictly positive because its nonnegative
weights sum to one. Taking logarithms of the exact tensor-power identity
multiplies the centered one-copy logarithm by $k$, proving the first bound.
The actual tensor-power state is positive, has trace one, and is invariant
under copy permutations. The quantitative signed-moment comparison therefore
gives the second bound, with the indicated common interval.
\end{proof}
